% ===========================================================================
% Section: System Overview
% ===========================================================================
\section{System Overview}
\label{sec:system_overview}

\winlic{} is designed as a modular, offline-first reconstruction framework
that processes pre-recorded multi-modal data to produce a photo-realistic
3D Gaussian map. The system architecture follows three design principles:

\begin{itemize}
    \item \textbf{ROS-free:} All communication is via direct function calls
    and shared memory structures, not ROS topics or message serialization.
    This enables Windows deployment and eliminates the overhead of
    message-passing middleware.

    \item \textbf{Pluggable:} Sensor data sources are abstracted behind the
    \texttt{ISensorSource} interface, allowing offline file readers and
    live camera SDKs to be swapped without modifying the reconstruction
    pipeline.

    \item \textbf{Algorithm-preserving:} The Gaussian Splatting mapping
    algorithm (initialization, extension, optimization, densification) is
    preserved verbatim from \gaussianlic, ensuring that reconstruction
    quality is identical to the original system.
\end{itemize}

\subsection{Pipeline Architecture}

Figure~\ref{fig:pipeline} illustrates the full processing pipeline.

% ===========================================================================
% Figure: Full system pipeline overview
% Data sources → ISensorSource → DataAligner → DepthFusion → GaussianModel → PLY
% ===========================================================================
\begin{figure*}[t]
\centering
\begin{tikzpicture}[node distance=0.8cm and 1.0cm]

% ---- Row 1: Sensor sources ----
\node[sensor] (lidar)  {LiDAR\\PointCloud};
\node[sensor, right=of lidar] (itof)   {iToF\\Depth};
\node[sensor, right=of itof]  (stereo) {Stereo\\RGB};
\node[sensor, right=of stereo] (imu)   {IMU};

% ---- Row 2: ISensorSource interface ----
\node[block, below=1.2cm of itof, minimum width=9cm, fill=blue!12]
  (iface) {\texttt{ISensorSource} (abstract interface)};

% ---- Row 3: Concrete implementations (side by side) ----
\node[blockgreen, below=0.8cm of iface, xshift=-2.8cm, minimum width=3.5cm]
  (reader) {DatasetReader\\(offline files)};
\node[blockorange, below=0.8cm of iface, xshift=2.8cm, minimum width=3.5cm]
  (orbbec) {OrbbecSource\\(live SDK)};

% ---- Row 4: Data aligner ----
\node[block, below=of iface, yshift=-2.2cm, minimum width=4cm]
  (aligner) {DataAligner\\(timestamp sync)};

% ---- Row 5: Depth fusion ----
\node[blockred, below=of aligner, minimum width=4cm]
  (fusion) {Depth Fusion\\(LiDAR + iToF + Stereo)};

% ---- Row 6: Depth completion ----
\node[blockorange, below=of fusion, minimum width=4cm]
  (completer) {SPNet Depth\\Completion (TensorRT)};

% ---- Row 7: Gaussian model ----
\node[blockgreen, below=of completer, minimum width=4cm]
  (gaussian) {GaussianModel\\(3DGS)};

% ---- Row 8: Output ----
\node[block, below=of gaussian, minimum width=4cm, fill=gray!10]
  (output) {PLY Map +\\Metrics};

% ---- Arrows: sensors → interface ----
\draw[arrow] (lidar.south) -- (iface.north -| lidar.south);
\draw[arrow] (itof.south)  -- (iface.north -| itof.south);
\draw[arrow] (stereo.south)-- (iface.north -| stereo.south);
\draw[arrow] (imu.south)   -- (iface.north -| imu.south);

% ---- Arrows: interface → implementations ----
\draw[arrow] (iface.south -| reader.north) -- (reader.north);
\draw[arrow] (iface.south -| orbbec.north) -- (orbbec.north);

% ---- Arrows: implementations → aligner ----
\draw[arrow] (reader.south) |- (aligner.west);
\draw[arrow] (orbbec.south) |- (aligner.east);

% ---- Main pipeline arrows ----
\draw[arrow] (aligner) -- (fusion);
\draw[arrow] (fusion)  -- (completer);
\draw[arrow] (completer) -- (gaussian);
\draw[arrow] (gaussian) -- (output);

% ---- Group: Pluggable layer ----
\begin{scope}[on background layer]
  \node[group=Pluggable Sensor Layer, fit=(iface)(reader)(orbbec),
        fill=blue!3] (pluggable) {};
\end{scope}

\end{tikzpicture}
\caption{Overview of the \winlic{} processing pipeline. Multi-modal sensor
data enters through the pluggable \texttt{ISensorSource} abstraction, is
synchronized by the DataAligner, fused into a dense depth map, completed by
the SPNet TensorRT engine, and fed into the 3D Gaussian Splatting model for
reconstruction.}
\label{fig:pipeline}
\end{figure*}

The system comprises six stages:

\begin{enumerate}
    \item \textbf{Sensor Input.} Multi-modal data enters through an
    \texttt{ISensorSource} implementation---either a
    \texttt{DatasetReader} for offline files or an \texttt{OrbbecSource}
    for live camera streams.

    \item \textbf{Data Alignment.} The \texttt{DataAligner} synchronizes
    frames across modalities by timestamp, producing a single
    \texttt{AlignedFrame} struct containing the RGB image, LiDAR point
    cloud, depth maps from each source, IMU data, and camera pose.

    \item \textbf{Depth Fusion.} The \texttt{DepthFusion} module combines
    sparse LiDAR depth, dense iToF depth, and stereo SGM depth into a
    single confidence-weighted depth map (Section~\ref{sec:sensor_fusion}).

    \item \textbf{Depth Completion.} The fused depth map is completed by
    the SPNet TensorRT engine, which fills holes using RGB guidance.

    \item \textbf{Gaussian Reconstruction.} The completed depth map and
    fused point cloud are fed into the \texttt{GaussianModel}, which
    initializes, extends, and optimizes the 3D Gaussian map
    (Section~\ref{sec:gaussian_reconstruction}).

    \item \textbf{Output.} The final Gaussian map is exported as a PLY
    file, along with rendered images, depth visualizations, and quality
    metrics (PSNR, SSIM, LPIPS).
\end{enumerate}

\subsection{Data Structures}

The original \gaussianlic{} used ROS message types (\texttt{sensor\_msgs::Image},
\texttt{sensor\_msgs::PointCloud2}, \texttt{geometry\_msgs::PoseStamped})
as the inter-module data carriers. In \winlic, these are replaced by
ROS-free structs:

\begin{itemize}
    \item \texttt{PointXYZRGB}: A 3D point with 8-bit colour
    ($x, y, z, r, g, b$).
    \item \texttt{CameraPose}: Timestamped pose as an Eigen quaternion
    and translation vector.
    \item \texttt{ImuMeasurement}: Timestamped accelerometer and gyroscope
    readings.
    \item \texttt{AlignedFrame}: A self-contained frame containing the
    RGB image (\texttt{cv::Mat}), multiple depth maps, point cloud,
    pose, and metadata.
\end{itemize}

These structs use only OpenCV (\texttt{cv::Mat}) and Eigen
(\texttt{Eigen::Quaterniond}, \texttt{Eigen::Vector3d}) types, which are
fully portable across Windows and Linux.

\subsection{Build System}

The original system used \texttt{catkin} (ROS build system) with
\texttt{find\_package(catkin REQUIRED COMPONENTS roscpp rospy ...)}.
\winlic{} replaces this with a pure CMake build that links against
OpenCV, LibTorch (PyTorch C++ API), TensorRT, yaml-cpp, and CUDA.
The build is configured via CMake cache variables and supports both
MSVC (Windows) and GCC (Linux) out of the box.
